% block: shifted-nonhermitian-heading page:188 type:section
\subsection{{Shifted Non-Hermitian Linear Systems}}\label{sec:shifted-nonhermitian-systems}
{Throughout this section, the coefficient matrix $A^{(k)}$ in \eqref{eq:shifted-linear-systems} is assumed to be non-Hermitian, i.e., $A^{(k)}\ne(A^{(k)})^{\mathrm H}$.}
% block: shifted-bicg-heading page:188 type:subsection
\subsubsection{{The Shifted BiCG Method}}\label{subsec:shifted-bicg}
{As seen in Sect.\ \ref{subsec:shifted-cg}, the $i$th residual vector $\bm r_i$ of the CG method for the seed system $A\bm x=\bm b$ and the $i$th residual vector $\bm r_i^{\sigma}$ for the shifted system $(A+\sigma I)\bm x^{\sigma}=\bm b$ are collinear. Theorem\ \ref{thm:shifted-residual-collinearity} is a generalized result of the condition that two residual vectors of a Krylov subspace method are collinear.}
% block: shifted-collinearity-theorem pages:188,189 type:theorem
\par\Needspace{14\baselineskip}
\begin{Theorem}[\BookCite{67}]\label{thm:shifted-residual-collinearity}
{Let $W_1\subseteq W_2\subseteq\cdots\subseteq W_k$ be a sequence of subspaces of $\mathbb C^N$ such that $\dim(W_i)=i$ and $W_i\cap\mathcal K_{i+1}(A,\bm b)^{\perp}=\{\bm0\}$ for $i=1,2,\ldots,k$. Let $\bm x_i\in\mathcal K_i(A,\bm b)^{\perp}$ be an approximate solution of $A\bm x=\bm b$ defined via the following Petrov--Galerkin condition of the residual $\bm r_i=\bm b-A\bm x_i=p_i(A)\bm b$:}
\[\bm r_i\perp W_i\quad\text{for }i=1,2,\ldots,k,\]
{where $p_i(A)=\sum_{k=0}^i c_kA^i$. Similarly, let $\bm x_i^{\sigma}\in\mathcal K_i(A+\sigma I,\bm b)=\mathcal K_i(A,\bm b)$ be the approximation to the solution of $(A+\sigma I)\bm x^{\sigma}=\bm b$ with the residual $\bm r^{\sigma}=\bm b-(A+\sigma I)\bm x_i^{\sigma}=p_i^{\sigma}(A+\sigma I)\bm b$, again satisfying}
\[\bm r_i^{\sigma}\perp W_i\quad\text{for }i=1,2,\ldots,k.\]
{Then $\bm r_i$ and $\bm r_i^{\sigma}$ are collinear, i.e., there exists $\pi_i^{\sigma}\in\mathbb C$ such that $\bm r_i=\pi_i^{\sigma}\bm r_i^{\sigma}$.}
\end{Theorem}
% block: shifted-collinearity-proof page:189 type:proof
\begin{Proof}
{Let $U_1,U_2$ be subspaces of $\mathbb C^N$. Then from standard linear algebra, it follows that $(U_1\cap U_2)^{\perp}=U_1^{\perp}+U_2^{\perp}$. From assumption $W_i\cap\mathcal K_{i+1}(A,\bm b)^{\perp}=\{\bm0\}$, we have}
\begin{equation}\label{eq:shifted-collinearity-complement-sum}(W_i\cap\mathcal K_{i+1}(A,\bm b)^{\perp})^{\perp}=\{\bm0\}^{\perp}\Leftrightarrow W_i^{\perp}+\mathcal K_{i+1}(A,\bm b)=\mathbb C^N.\end{equation}
{$\dim W_i^{\perp}=N-i$ and $\mathcal K_{i+1}(A,\bm b)=i+1$, and thus $\dim W_i^{\perp}+\dim\mathcal K_{i+1}(A,\bm b)=N+1$. On the other hand, from \eqref{eq:shifted-collinearity-complement-sum}, $\dim(W_i^{\perp}+\mathcal K_{i+1}(A,\bm b))=\dim\mathbb C^N=N$. Thus}
\[\dim(W_i^{\perp}\cap\mathcal K_{i+1}(A,\bm b))=1.\]
{Since $\bm r_i,\bm r_i^{\sigma}\in\mathcal K_{i+1}(A,\bm b)\cap W_i^{\perp}$, two residual vectors $\bm r_i,\bm r_i^{\sigma}$ lie in the same one-dimensional subspace, which means that $\bm r_i,\bm r_i^{\sigma}$ are collinear.}\par\nopagebreak[4]\end{Proof}
% block: shifted-bicg-application page:189 type:ordinary-paragraph
{Theorem\ \ref{thm:shifted-residual-collinearity} indicates that using the relation $W_n=\mathcal K_n(A^{\mathrm H},\bm r_0^*)$, the BiCG residual vector $\bm r_n$ from $A\bm x=\bm b$ and the BiCG residual vector $\bm r_n^{(k)}$ from $(A+\sigma_k I)\bm x^{(k)}=\bm b$ are collinear, i.e.,}
\[\bm r_n=\pi_n^{(k)}\bm r_n^{(k)}.\]
{Thus, following the derivation of the shifted CG method in Sect.\ \ref{subsec:shifted-cg}, the algorithm of the shifted BiCG method is obtained, which is listed in Algorithm\ \ref{alg:shifted-bicg}. Similarly, the BiCR method in Sect.\ \ref{subsec:bicr-method} is developed to solve shifted linear systems, which is referred to as the shifted BiCR method in \BookCite{85}.}
% block: shifted-bicgstab-heading page:189 type:subsection
\subsubsection{{The Shifted BiCGSTAB Method}}\label{subsec:shifted-bicgstab}
{The shifted BiCGSTAB method and the shifted BiCGSTAB($\ell$) method are proposed in \BookCite{67}. Since the derivation of the shifted BiCGSTAB($\ell$) method is somewhat complicated, the derivation of the shifted BiCGSTAB method is described, which is based on the explanation in \BookCite{140}. For the derivation of the shifted BiCGSTAB method, we consider the seed system $A\bm x=\bm b$ and the shifted system $(A+\sigma I)\bm x^{\sigma}=\bm b$. First of all, let us recall the BiCGSTAB method in Sect.\ \ref{subsec:product-type-krylov}. Let $\bm r_{n+1}$ be the residual vector of the BiCGSTAB method. Then $\bm r_{n+1}$ and the other iterates $\bm t_n,\bm p_n,\bm x_n$ are described by}
% block: shifted-bicg-algorithm page:190 type:ordinary-paragraph
\begin{algorithm}[H]\caption{{The shifted BiCG method}}\label{alg:shifted-bicg}\begin{algorithmic}[1]
\Statex\textbf{Input:} {Choose a seed system $s\in S=\{1,2,\ldots,m\}$ and set $\bm r_0^{(s)}=\bm b,\ \beta_{-1}^{(s)}=0$}
\Statex\textbf{Input:} $A$ {and} $\sigma_s$ {for} $k=1,2,\ldots,m$
\Statex\textbf{Input:} $\bm x_0^{(k)}=\bm p_{-1}^{(k)}=\bm p_{-1}^{*(k)}=\bm0,\ \pi_0^{(s,k)}=\pi_{-1}^{(s,k)}=1$ {for} $k=1,2,\ldots,m$
\Statex\textbf{Input:} {Choose $\bm r_0^*\in\mathbb C^N$, e.g., $\bm r_0^*=\bm r_0$}
\Statex\textbf{Output:} $\bm x_n^{(k)}$ {for} $k=1,2,\ldots,m$
\For{$n=0,1,\ldots$, {until convergence}}
\State $\bm p_n^{(s)}=\bm r_n^{(s)}+\beta_{n-1}^{(s)}\bm p_{n-1}^{(s)},\quad\bm p_n^{*(s)}=\bm r_n^{*(s)}+\bar\beta_{n-1}^{(s)}\bm p_{n-1}^{*(s)}$
\State $\alpha_n^{(s)}=\dfrac{(\bm r_n^{*(s)},\bm r_n^{(s)})}{(\bm p_n^{*(s)},A^{(s)}\bm p_n^{(s)})}$
\State $\bm x_{n+1}^{(s)}=\bm x_n^{(s)}+\alpha_n^{(s)}\bm p_n^{(s)}$
\State \{{Begin shifted system}\}
\State {Run lines 6--15 of the shifted CG method (Algorithm\ \ref{alg:shifted-cg}).}
\State \{{End shifted system}\}
\State $\bm r_{n+1}^{(s)}=\bm r_n^{(s)}-\alpha_n^{(s)}A^{(s)}\bm p_n^{(s)},\quad\bm r_{n+1}^{*(s)}=\bm r_n^{*(s)}-\bar\alpha_nA^{(s)\mathrm H}\bm p_n^{*(s)}$
\State $\beta_n^{(s)}=\dfrac{(\bm r_{n+1}^{*(s)},\bm r_{n+1}^{(s)})}{(\bm r_n^{*(s)},\bm r_n^{(s)})}$
\EndFor\end{algorithmic}\end{algorithm}
% block: bicgstab-polynomial-recollection page:190 type:ordinary-paragraph
\[\begin{aligned}\bm p_n&:=Q_n(A)\bm p_n^{\mathrm{BiCG}}=\bm r_n+\beta_{n-1}(\bm p_{n-1}-\zeta_{n-1}A\bm p_{n-1})\\&=\bm r_n+\beta_{n-1}\left[\bm p_{n-1}+\frac{\zeta_{n-1}}{\alpha_{n-1}}(\bm t_{n-1}-\bm r_{n-1})\right],\\
\bm t_n&:=Q_n(A)\bm r_{n+1}^{\mathrm{BiCG}}=\bm r_n-\alpha_nA\bm p_n,\\
\bm r_{n+1}&:=Q_{n+1}(A)\bm r_{n+1}^{\mathrm{BiCG}}=\bm t_n-\zeta_nA\bm t_n,\\
\bm x_{n+1}&:=\bm x_n+\alpha_n\bm p_n+\zeta_n\bm t_n,\end{aligned}\]
{where $\bm r_{n+1}^{\mathrm{BiCG}}$ and $\bm p_n^{\mathrm{BiCG}}$ are the BiCG residual vector at $n+1$ iteration step and the BiCG search direction at $n$ iteration step respectively, and the polynomial $Q_{n+1}(z)$ is defined by}
\begin{equation}\label{eq:shifted-bicgstab-seed-polynomial-initial}Q_0(\lambda):=1,\end{equation}
\begin{equation}\label{eq:shifted-bicgstab-seed-polynomial-recurrence}Q_{n+1}(\lambda):=(1-\zeta_n\lambda)Q_n(\lambda),\quad n=0,1,\ldots\end{equation}
{As for the BiCGSTAB method, $\zeta_n$ is determined so that $\|\bm r_{n+1}\|$ is minimized. Recall that the residual vector of the shifted BiCG method in Sect.\ \ref{subsec:shifted-bicg} for $(A+\sigma I)\bm x^{\sigma}=\bm b$ is written as}
\[\bm r_{n+1}^{\sigma,\mathrm{BiCG}}=\xi_{n+1}^{\sigma}\bm r_{n+1}^{\mathrm{BiCG}},\quad\xi_{n+1}^{\sigma}\in\mathbb C.\]
% block: shifted-bicgstab-xi-recurrence pages:190,191 type:ordinary-paragraph
{Here, from \eqref{eq:shifted-cg-pi-recurrence}, the scalar $\xi_{n+1}^{\sigma}$ is defined by $\xi_{n+1}^{\sigma}:=(\pi_{n+1}^{\sigma})^{-1}$ for all $n$, and thus we have}
\begin{equation}\label{eq:shifted-bicgstab-xi-recurrence}\begin{aligned}\xi_{n+1}^{\sigma}&=\frac1{\left(1+\frac{\beta_{n-1}}{\alpha_{n-1}}\alpha_n+\alpha_n\sigma\right)\pi_n^{\sigma}-\frac{\beta_{n-1}}{\alpha_{n-1}}\alpha_n\pi_{n-1}^{\sigma}}\\
&=\frac1{\left(1+\frac{\beta_{n-1}}{\alpha_{n-1}}\alpha_n+\alpha_n\sigma\right)(\xi_n^{\sigma})^{-1}-\frac{\beta_{n-1}}{\alpha_{n-1}}\alpha_n(\xi_{n-1}^{\sigma})^{-1}}\\
&=\frac{\xi_n^{\sigma}\xi_{n-1}^{\sigma}\alpha_{n-1}}{(1+\alpha_n\sigma)\xi_{n-1}^{\sigma}\alpha_{n-1}+\alpha_n\beta_{n-1}(\xi_{n-1}^{\sigma}-\xi_n^{\sigma})},\end{aligned}\end{equation}
{where $\xi_{-1}^{\sigma}=\xi_0^{\sigma}=1$.}
% block: shifted-bicgstab-smoothing page:191 type:ordinary-paragraph
{We are now ready to describe the shifted BiCGSTAB method. The residual vector of the shifted BiCGSTAB method is defined by}
\[\bm r_{n+1}^{\sigma,\mathrm{STAB}}:=Q_{n+1}^{\sigma}(A)\bm r_{n+1}^{\sigma,\mathrm{BiCG}},\]
{where $Q_{n+1}^{\sigma}(\lambda)$ are recursively defined as}
\begin{equation}\label{eq:shifted-bicgstab-shifted-polynomial}Q_{n+1}^{\sigma}(\lambda):=[1-\zeta_n^{\sigma}(\lambda+\sigma)]Q_n^{\sigma}(\lambda),\quad n=0,1,\ldots\end{equation}
{with $Q_0^{\sigma}(\lambda):=1$. Here, $\zeta_n^{\sigma}$ is determined so that the shifted BiCGSTAB residual and the BiCGSTAB residual are colinear, i.e.,}
\[\bm r_{n+1}^{\sigma,\mathrm{STAB}}=c_n\bm r_{n+1}^{\mathrm{STAB}},\quad c_n\in\mathbb C.\]
{To this end, $\zeta_n^{\sigma}$ is determined so that}
\begin{equation}\label{eq:shifted-bicgstab-polynomial-collinearity}Q_{n+1}^{\sigma}(\lambda)=\tau_{n+1}^{\sigma}Q_{n+1}(\lambda),\quad\tau_{n+1}^{\sigma}\in\mathbb C.\end{equation}
{Then, the residual vector of the shifted BiCGSTAB method is defined and written as}
\begin{equation}\label{eq:shifted-bicgstab-residual-collinearity}\begin{aligned}\bm r_{n+1}^{\sigma,\mathrm{STAB}}&:=Q_{n+1}^{\sigma}(A)\bm r_{n+1}^{\sigma,\mathrm{BiCG}}\\&=\tau_{n+1}^{\sigma}Q_{n+1}(A)\xi_{n+1}^{\sigma}\bm r_{n+1}^{\mathrm{BiCG}}\\&=\tau_{n+1}^{\sigma}\xi_{n+1}^{\sigma}Q_{n+1}(A)\bm r_{n+1}^{\mathrm{BiCG}}\\&=\tau_{n+1}^{\sigma}\xi_{n+1}^{\sigma}\bm r_{n+1}^{\mathrm{STAB}}.\end{aligned}\end{equation}
% block: shifted-bicgstab-tau-zeta pages:191,192 type:ordinary-paragraph
{In what follows, the parameters $\tau_{n+1}^{\sigma},\xi_{n+1}^{\sigma}$ in \eqref{eq:shifted-bicgstab-residual-collinearity} are determined. It follows from \eqref{eq:shifted-bicgstab-seed-polynomial-recurrence}, \eqref{eq:shifted-bicgstab-shifted-polynomial}, and \eqref{eq:shifted-bicgstab-polynomial-collinearity} that}
\begin{equation}\label{eq:shifted-bicgstab-polynomial-matching}\underbrace{\tau_{n+1}^{\sigma}(1-\zeta_n\lambda)Q_n(\lambda)}_{Q_{n+1}^{\sigma}(\lambda)}=[1-\zeta_n^{\sigma}(\lambda+\sigma)]\underbrace{\tau_n^{\sigma}Q_n(\lambda)}_{Q_n^{\sigma}(\lambda)}.\end{equation}
{Comparing the coefficients on both sides of \eqref{eq:shifted-bicgstab-polynomial-matching} yields}
\[\zeta_n\tau_{n+1}^{\sigma}=\zeta_n^{\sigma}\tau_n^{\sigma},\quad\tau_{n+1}^{\sigma}=(1-\zeta_n^{\sigma}\sigma)\tau_n^{\sigma},\]
{from which, we obtain}
\begin{equation}\label{eq:shifted-bicgstab-tau-zeta}\tau_{n+1}^{\sigma}=\frac{\tau_n^{\sigma}}{1+\zeta_n\sigma},\quad\zeta_n^{\sigma}=\frac{\zeta_n}{1+\zeta_n\sigma},\end{equation}
{where $\tau_0^{\sigma}=1$ from the relation \eqref{eq:shifted-bicgstab-polynomial-collinearity} and $Q_0^{\sigma}(\lambda)=Q_0(\lambda)=1$. The shifted BiCGSTAB iterates for $(A+\sigma I)\bm x^{\sigma}=\bm b$ are now given by}
\begin{equation}\label{eq:shifted-bicgstab-direction}\begin{aligned}\bm p_n^{\sigma}&:=Q_n^{\sigma}(A)\bm p_n^{\sigma,\mathrm{BiCG}}=\bm r_n^{\sigma}+\beta_{n-1}^{\sigma}\left[\bm p_{n-1}^{\sigma}+\frac{\zeta_{n-1}^{\sigma}}{\alpha_{n-1}^{\sigma}}(\bm t_{n-1}^{\sigma}-\bm r_{n-1}^{\sigma})\right]\\&=\tau_n^{\sigma}\xi_n^{\sigma}\bm r_n+\beta_{n-1}^{\sigma}\left[\bm p_{n-1}^{\sigma}+\frac{\zeta_{n-1}^{\sigma}}{\alpha_{n-1}^{\sigma}}\tau_{n-1}^{\sigma}(\xi_n^{\sigma}\bm t_{n-1}-\xi_{n-1}^{\sigma}\bm r_{n-1})\right],\end{aligned}\end{equation}
\begin{equation}\label{eq:shifted-bicgstab-intermediate}\bm t_n^{\sigma}:=Q_n^{\sigma}(A)\bm r_{n+1}^{\sigma,\mathrm{BiCG}}=\tau_n^{\sigma}\xi_{n+1}^{\sigma}Q_n\bm r_{n+1}^{\mathrm{BiCG}}=\tau_n^{\sigma}\xi_{n+1}^{\sigma}\bm t_n,\end{equation}
\begin{equation}\label{eq:shifted-bicgstab-residual-update}\bm r_{n+1}^{\sigma}:=Q_{n+1}^{\sigma}(A)\bm r_{n+1}^{\sigma,\mathrm{BiCG}}=\tau_{n+1}^{\sigma}\xi_{n+1}^{\sigma}Q_{n+1}\bm r_{n+1}^{\mathrm{BiCG}}=\tau_{n+1}^{\sigma}\xi_{n+1}^{\sigma}\bm r_{n+1}^{\mathrm{STAB}},\end{equation}
\begin{equation}\label{eq:shifted-bicgstab-solution-update}\bm x_{n+1}^{\sigma}:=\bm x_n^{\sigma}+\alpha_n^{\sigma}\bm p_n^{\sigma}+\zeta_n^{\sigma}\bm t_n^{\sigma}.\end{equation}
{From \eqref{eq:shifted-bicgstab-xi-recurrence}, \eqref{eq:shifted-bicgstab-tau-zeta}, \eqref{eq:shifted-bicgstab-direction}--\eqref{eq:shifted-bicgstab-solution-update}, the shifted BiCGSTAB method is obtained, and the algorithm is written in Algorithm\ \ref{alg:shifted-bicgstab}.}
% block: shifted-gmres-heading page:192 type:subsection
\subsubsection{{The Shifted GMRES Method}}\label{subsec:shifted-gmres}
{In this subsection, the shifted GMRES method \BookCite{43} is derived. To this end, we consider applying the GMRES method to the shifted linear systems $(A+\sigma I)\bm x=\bm b$. For simplicity, the initial guess is set to $\bm x_0=\bm0$. Then the GMRES method finds an approximate solution over the following Krylov subspace:}
\begin{equation}\label{eq:shifted-gmres-krylov-membership}\bm x_n^{\sigma}\in\mathcal K_n(A+\sigma I,\bm b).\end{equation}
{It follows from the shift-invariance property \eqref{eq:krylov-shift-invariance} that we have $\mathcal K_n(A+\sigma I,\bm b)=\mathcal K_n(A,\bm b)$. Thus when using the Arnoldi process, the orthonormal basis vectors of $\mathcal K_n(A+\sigma I,\bm b)$ are equivalent to those of $\mathcal K_n(A,\bm b)$, i.e., $V_n$ in Sect.\ \ref{subsec:arnoldi-process}. Then, \eqref{eq:shifted-gmres-krylov-membership} and the corresponding residual are written as}
\[\bm x_n^{\sigma}=V_n\bm y_n,\quad\bm y_n\in\mathbb C^n\]
% block: shifted-bicgstab-algorithm page:193 type:ordinary-paragraph
\begin{algorithm}[H]\caption{{The shifted BiCGSTAB method}}\label{alg:shifted-bicgstab}\begin{algorithmic}[1]
\Statex\textbf{Input:} $\bm x_0=\bm0,\ \bm p_{-1}=\bm0,\ \bm r_0=\bm b,\ \beta_{-1}=0$
\Statex\textbf{Input:} {Choose $\bm r_0^*\in\mathbb C^N$, e.g., $\bm r_0^*=\bm r_0$}
\Statex\textbf{Input:} $\bm x_0^{(k)}=\bm0,\ \bm p_0^{(k)}=\bm b,\ \beta_{-1}^{(k)}=0,\ \alpha_{-1}^{(k)}=\xi_{-1}^{(k)}=\xi_0^{(k)}=\tau_0^{(k)}=1$ {for} $k=1,2,\ldots,m$
\Statex\textbf{Output:} $\bm x_n^{(k)}$
\For{$n=0,1,\ldots$}
\State $\bm p_n=\bm r_n+\beta_{n-1}(\bm p_{n-1}-\zeta_{n-1}A\bm p_{n-1})$
\State $\alpha_n=\dfrac{(\bm r_0^*,\bm r_n)}{(\bm r_0^*,A\bm p_n)}$
\State $\bm t_n=\bm r_n-\alpha_nA\bm p_n$
\State $\zeta_n=\dfrac{(A\bm t_n,\bm t_n)}{(A\bm t_n,A\bm t_n)}$
\State $\bm x_{n+1}=\bm x_n+\alpha_n\bm p_n+\zeta_n\bm t_n$
\State $\bm r_{n+1}=\bm t_n-\zeta_nA\bm t_n$
\State $\beta_n=\dfrac{\alpha_n}{\zeta_n}\times\dfrac{(\bm r_0^*,\bm r_{n+1})}{(\bm r_0^*,\bm r_n)}$
\State \{{Begin shifted system}\}
\For{$k=1,2,\ldots,m$}
\State $\xi_{n+1}^{(k)}=\dfrac{\xi_n^{(k)}\xi_{n-1}^{(k)}\alpha_{n-1}}{(1+\alpha_n\sigma)\xi_{n-1}^{(k)}\alpha_{n-1}+\alpha_n\beta_{n-1}(\xi_{n-1}^{(k)}-\xi_n^{(k)})}$
\State $\alpha_n^{(k)}=\dfrac{\xi_{n+1}^{(k)}}{\xi_n^{(k)}}\alpha_k$
\State $\zeta_n^{(k)}=\dfrac{\zeta_n}{1+\zeta_n\sigma}$
\State $\bm x_{n+1}^{(k)}=\bm x_n^{(k)}+\alpha_n^{(k)}\bm p_n^{(k)}+\zeta_n^{(k)}\bm t_n^{(k)}$
\State $\tau_{n+1}^{(k)}=\dfrac{\tau_n^{(k)}}{1+\zeta_n\sigma}$
\State $\beta_n^{(k)}=\left(\dfrac{\xi_{n+1}^{(k)}}{\xi_n^{(k)}}\right)^2\beta_k$
\State $\begin{aligned}\bm p_{n+1}^{(k)}&=\tau_{n+1}^{(k)}\xi_{n+1}^{(k)}\bm r_{n+1}\\&\quad+\beta_n^{(k)}\left[\bm p_n^{(k)}+\frac{\zeta_n^{(k)}}{\alpha_n^{(k)}}\tau_n^{(k)}(\xi_{n+1}^{(k)}\bm t_n^{(k)}-\xi_n^{(k)}\bm r_n)\right]\end{aligned}$
\EndFor
\State \{{End shifted system}\}
\EndFor\end{algorithmic}\end{algorithm}
% block: shifted-gmres-residual-expansion page:193 type:ordinary-paragraph
{and}
\[\begin{aligned}\bm r_n^{\sigma}&=\bm b-(A+\sigma I)\bm x_n^{\sigma}\\
&=\|\bm b\|V_{n+1}\bm e_1-AV_n\bm y_n-\sigma V_n\bm y_n\\
&=\|\bm b\|V_{n+1}\bm e_1-V_{n+1}H_{n+1,n}\bm y_n-\sigma V_{n+1}\begin{bmatrix}\bm y_n\\0\end{bmatrix}\\
&=\|\bm b\|V_{n+1}\bm e_1-V_{n+1}H_{n+1,n}\bm y_n-\sigma V_{n+1}\begin{bmatrix}I_n\\\bm0^{\top}\end{bmatrix}\bm y_n\\
&=V_{n+1}\left[\|\bm b\|\bm e_1-\left(H_{n+1,n}+\begin{bmatrix}\sigma I_n\\\bm0^{\top}\end{bmatrix}\right)\bm y_n\right]\\
&=V_{n+1}(\|\bm b\|\bm e_1-H_{n+1,n}^{\sigma}\bm y_n),\end{aligned}\]
% block: shifted-gmres-least-squares page:194 type:ordinary-paragraph
{where $I_n$ is the $n\times n$ identity matrix and $H_{n+1,n}^{\sigma}$ is a so-called shifted Hessenberg matrix whose elements are $(H_{n+1,n}^{\sigma})_{i,j}=(H_{n+1,n})_{i,j}$ for $i\ne j$ and $(H_{n+1,n}^{\sigma})_{i,j}=(H_{n+1,n})_{i,j}+\sigma$ for $i=j$. Then similar to \eqref{eq:gmres-projected-minimization}, the residual norm can be minimized by solving}
\[\bm y_n^{\sigma}:=\operatorname*{arg\,min}_{\bm y\in\mathbb C^n}\|\beta\bm e_1-H_{n+1,n}^{\sigma}\bm y\|.\]
{The above least-squares problems can be efficiently solved by Givens rotations that are described in Sect.\ \ref{subsec:minres}. Using $\bm y_n^{\sigma}$, the approximate solution is now given by}
\[\bm x_n^{\sigma}=V_n\bm y_n^{\sigma}.\]
{If the coefficient matrix $A$ is Hermitian and $\sigma\in\mathbb R$, then the shifted GMRES method reduces to the shifted MINRES method in Sect.\ \ref{subsec:shifted-minres}. The algorithm of the shifted GMRES method is written in Algorithm\ \ref{alg:shifted-gmres}.}
% block: shifted-gmres-algorithm page:194 type:ordinary-paragraph
\begin{algorithm}[H]\caption{{The shifted GMRES method}}\label{alg:shifted-gmres}\begin{algorithmic}[1]
\Statex\textbf{Input:} $\sigma_k\ (k=1,2,\ldots,m)$
\Statex\textbf{Output:} $\bm x_n^{(k)}\ (k=1,2,\ldots,m)$
\State $\bm g=(\|\bm b\|,0,\ldots,0)^{\top},\ \bm v_1=\bm b/\|\bm b\|$
\For{$n=1,2,\ldots$}
\State ({Arnoldi process})
\State $\bm t=A\bm v_n$
\For{$i=1,2,\ldots,n$}
\State $h_{i,n}=(\bm v_i,\bm t)$
\State $\bm t=\bm t-h_{i,n}\bm v_i$
\EndFor
\State $h_{n+1,n}=\|\bm t\|$
\State $\bm v_{n+1}=\bm t/h_{n+1,n}$
\State \{{Begin shifted system}\}
\State $h_{i,n}^{(k)}=h_{i,n}$ {for} $i\ne n$
\State $h_{n,n}^{(k)}=h_{n,n}+\sigma_k$
\For{$k=1,2,\ldots,m$}
\State \{{Givens rotations}\}
\For{$i=1,2,\ldots,n-1$}
\State $\begin{bmatrix}h_{i,n}^{(k)}\\h_{i+1,n}^{(k)}\end{bmatrix}=\begin{bmatrix}c_i^{(k)}&s_i^{(k)}\\-\bar s_i^{(k)}&c_i^{(k)}\end{bmatrix}\begin{bmatrix}h_{i,n}^{(k)}\\h_{i+1,n}^{(k)}\end{bmatrix}$
\EndFor
\State $c_n^{(k)}=\dfrac{|h_{n,n}^{(k)}|}{\sqrt{|h_{n,n}^{(k)}|^2+|h_{n+1,n}^{(k)}|^2}}$
\State $\bar s_n=\dfrac{h_{n+1,n}^{(k)}}{h_{n,n}^{(k)}}c_n$
\State $h_{n,n}^{(k)}=c_n^{(k)}h_{n,n}^{(k)}+s_n^{(k)}h_{n+1,n}^{(k)}$
\State $h_{n+1,n}^{(k)}=0$
\State $\begin{bmatrix}g_n^{(k)}\\g_{n+1}^{(k)}\end{bmatrix}=\begin{bmatrix}c_n^{(k)}&s_n^{(k)}\\-\bar s_n^{(k)}&c_n^{(k)}\end{bmatrix}\begin{bmatrix}g_n^{(k)}\\0\end{bmatrix}$
\State ({Check convergence})
\If{$|g_{n+1}^{(k)}|/\|\bm b\|\le\epsilon$,}
\State $\bm x_n^{(k)}=V_n(H_n^{(k)})^{-1}\bm g^{(k)}$
\EndIf\EndFor
\State \{{End shifted system}\}
\EndFor\end{algorithmic}\end{algorithm}
% block: shifted-gmres-restart-difficulty pages:194,195 type:ordinary-paragraph
{The differences between the GMRES method and the shifted GMRES method are the initial guess $\bm x_0=\bm0$ and line 13 in Algorithm\ \ref{alg:shifted-gmres} which corresponds to producing $H_{n+1,n}^{\sigma}$. Note that after Givens rotations $H_n$ becomes the upper triangular matrix, and thus $\bm c^{(k)}:=(H_n^{(k)})^{-1}\bm g$ in line 26 can easily be obtained by solving $H_n^{(k)}\bm c^{(k)}=\bm g$.}
{As mentioned in Sect.\ \ref{subsec:gmres-method}, the GMRES method is not practical when the number of iterations is large, due to growing memory requirement and computational costs. Instead, the restarted GMRES method is useful in practice, and was described in Algorithm\ \ref{alg:gmres-restarted}. For the same reason, the shifted GMRES method is not practical when the number of iterations is large. It is therefore natural to consider the restarted version for the shifted linear systems. However, if we use the idea of Algorithm\ \ref{alg:gmres-restarted}, then we face a difficulty in that after the restart the initial residual vectors $\bm r_0^{(i)}$ and $\bm r_0^{(j)}$ are, in general, not colinear any longer, and thus we cannot share the basis vectors of Krylov subspaces for solving shifted linear systems.}
% block: restarted-shifted-gmres-collinearity page:195 type:ordinary-paragraph
{Frommer and Gräsner \BookCite{68} nicely circumvented the difficulty. In what follows, the idea and the corresponding algorithm are described. Let $\bm r_n$ be the residual vector of the GMRES method for $A\bm x=\bm b$. Then they consider forcing the residual vector $\bm r_n^{\sigma}$ for the shifted system $(A+\sigma I)\bm x^{\sigma}=\bm b$ to become collinear with $\bm r_n$, i.e., find the approximate solution $\bm x^{\sigma}=\bm x_0^{\sigma}+V_n\bm y_n^{\sigma}$ such that}
\[\bm r_n^{\sigma}=\beta_n\bm r_n.\]
{$\bm r_n$ is written as}
\[\bm r_n=\bm b-A\bm x_n=V_{n+1}\bm z_{n+1},\]
{where $\bm z_{n+1}:=\|\bm r_0\|\bm e_1-H_{n+1,n}\bm y_n$. We now have the following equation:}
\[\begin{aligned}\bm r_n^{\sigma}=\beta_n\bm r_n&\Leftrightarrow\bm b-A^{\sigma}\bm x_n^{\sigma}=\beta_nV_{n+1}\bm z_{n+1}\\
&\Leftrightarrow\bm b-A^{\sigma}(\bm x_0^{\sigma}+V_n\bm y_n^{\sigma})=\beta_nV_{n+1}\bm z_{n+1}\\
&\Leftrightarrow\beta_0\bm r_0-A^{\sigma}V_n\bm y_n^{\sigma}=\beta_nV_{n+1}\bm z_{n+1}\\
&\Leftrightarrow\beta_0\bm r_0-V_{n+1}H_{n+1,n}^{\sigma}\bm y_n^{\sigma}=\beta_nV_{n+1}\bm z_{n+1}\\
&\Leftrightarrow\beta_0\bm r_0=V_{n+1}(\beta_n\bm z_{n+1}+H_{n+1,n}^{\sigma}\bm y_n^{\sigma})\\
&\Leftrightarrow V_{n+1}(H_{n+1,n}^{\sigma}\bm y_n^{\sigma}+\beta_n\bm z_{n+1})=\beta_0\bm r_0\\
&\Leftrightarrow V_{n+1}(H_{n+1,n}^{\sigma}\bm y_n^{\sigma}+\beta_n\bm z_{n+1})=\beta_0\|\bm r_0\|V_{n+1}\bm e_1\\
&\Leftrightarrow H_{n+1,n}^{\sigma}\bm y_n^{\sigma}+\beta_n\bm z_{n+1}=\beta_0\|\bm r_0\|\bm e_1,\end{aligned}\]
{from which $\beta_n$ and $\bm y_n^{\sigma}$ are determined by solving the following $(n+1)\times(n+1)$ linear systems:}
\begin{equation}\label{eq:restarted-shifted-gmres-collinearity-system}[H_{n+1,n}^{\sigma}\ \bm z_{n+1}]\begin{bmatrix}\bm y_n^{\sigma}\\\beta_n\end{bmatrix}=\beta_0\|\bm r_0\|\bm e_1.\end{equation}
% block: restarted-shifted-gmres-monitoring pages:195,196 type:ordinary-paragraph
{After the restart, the new initial residual vectors $\bm r_0(=\bm r_n),\bm r_0^{\sigma}(=\bm r_n^{\sigma})$ become collinear. Thus we can use Krylov subspace $\mathcal K_n(A,\bm b)$ for $A\bm x=\bm b$ to solve shifted linear systems $(A+\sigma I)\bm x^{\sigma}=\bm b^{\sigma}$. The algorithm of the restarted shifted GMRES method is given in Algorithm\ \ref{alg:restarted-shifted-gmres}.}
{Note that from the relation $\|\bm r_n^{(k)}\|=\|\beta_n^{(k)}\bm r_n\|$ one can check the convergence of $\|\bm r_n^{(k)}\|$ of Algorithm\ \ref{alg:restarted-shifted-gmres} by monitoring $|\beta_n^{(k)}|\times\|\bm r_n\|$. For further developments of the (restarted) shifted GMRES method, see \BookCite{55}, \BookCite{105}, \BookCite{115}, \BookCite{175}, and the references therein.}
% block: shifted-fom-initial-observation page:196 type:ordinary-paragraph
{The full orthogonalization method (FOM) \BookCite{149} is an extension of the CG method to solving non-Hermitian linear systems. As well as the CG method, the residual vector $\bm r_n^{\mathrm{FOM}}$ of the FOM satisfies}
\[\bm r_n^{\mathrm{FOM}}(\in\mathcal K_{n+1}(A,\bm b))\perp\mathcal K_n(A,\bm b).\]
% block: restarted-shifted-gmres-algorithm page:196 type:ordinary-paragraph
\begin{algorithm}[H]\caption{{The restarted shifted GMRES method}}\label{alg:restarted-shifted-gmres}\begin{algorithmic}[1]
\Statex\textbf{Input:} $\sigma_k\ (k=1,2,\ldots,m)$
\Statex\textbf{Input:} $\bm x_0\in\mathbb C^N$
\Statex\textbf{Input:} {Set $\bm x_0^{(k)}\in\mathbb C^N$ such that $\bm r_0^{(k)}=\beta_0^{(k)}\bm r_0$, e.g., $\bm x_0=\bm x_0^{(k)}=\bm0$ for all $k$.}
\Statex\textbf{Output:} $\bm x_n^{(k)}\ (k=1,2,\ldots,m)$
\State $\bm r_0=\bm b-A\bm x_0,\ \beta=\|\bm r_0\|$
\State {Run the Arnoldi process in Algorithm\ \ref{alg:arnoldi-modified} with $\bm v_1=\bm r_0/\beta$.}
\State {Compute $\bm y_n$ such that $\|\beta\bm e_1-H_{n+1,n}\bm y_n\|$ is minimized.}
\State $\bm x_n=\bm x_0+V_n\bm y_n,\ \bm z_{n+1}=\beta\bm e_1-H_{n+1,n}\bm y_n$
\State \{{Shifted systems}\}
\For{$k=1,2,\ldots,m$}
\State $H_{n+1,n}^{(k)}=H_{n+1,n}+\begin{bmatrix}\sigma_k I_n\\\bm0^{\top}\end{bmatrix}$
\State {Solve} $[H_{n+1,n}^{(k)}\ \bm z_{n+1}]\begin{bmatrix}\bm y_n^{(k)}\\\beta_n^{(k)}\end{bmatrix}=\beta_0^{(k)}\beta\bm e_1$.
\State $\bm x_n^{(k)}=\bm x_0^{(k)}+V_m\bm y_m^{(k)}$
\EndFor
\If{{not convergence}}
\State {Set $\bm x_0=\bm x_n,\ \bm x_0^{(k)}=\bm x_n^{(k)},\ \beta_0^{(k)}=\beta_n^{(k)}$, and go to line 1.}
\EndIf\end{algorithmic}\end{algorithm}
% block: restarted-shifted-fom pages:196,197 type:ordinary-paragraph
{When the FOM is applied to shifted linear systems, we have}
\[\bm r_n^{\sigma,\mathrm{FOM}}(\in\mathcal K_{n+1}(A+\sigma I,\bm b))\perp\mathcal K_n(A+\sigma I,\bm b).\]
{Then from the shift invariance property $\mathcal K_n(A+\sigma I,\bm b)=\mathcal K_n(A,\bm b)$, it follows that}
\[\bm r_n^{\mathrm{FOM}},\bm r_n^{\sigma,\mathrm{FOM}}\in\mathcal K_{n+1}(A,\bm b)\cap\mathcal K_n(A,\bm b)^{\perp}.\]
{From \eqref{eq:shifted-cg-one-dimensional-intersection}, the two residual vectors $\bm r_n^{\mathrm{FOM}},\bm r_n^{\sigma,\mathrm{FOM}}$ belong to a one-dimensional subspace. Thus the residual vectors applied to $A\bm x=\bm b$ and $(A+\sigma I)\bm x^{\sigma}=\bm b$ are collinear. This implies that we do not need to consider the trick \eqref{eq:restarted-shifted-gmres-collinearity-system} for keeping collinearity when considering the restart, since the residual vectors are collinear at each iteration step. For the details of the restarted shifted FOM, see \BookCite{158}. Further development of the restarted shifted FOM, see \BookCite{55}, \BookCite{112}, and the references therein.}
% block: shifted-idr-heading page:197 type:subsection
\subsubsection{{The Shifted IDR($s$) Method}}\label{subsec:shifted-idr}
{The IDR($s$) method in Sect.\ \ref{subsec:idr-method} is extended to solving shifted linear systems and the algorithm is referred to as the shifted IDR($s$) method \BookCite{48}. Prior to the derivation of the shifted IDR($s$) method, a variant of the IDR theorem (Theorem\ \ref{thm:idr-dimension-reduction}) is described next.}
% block: shifted-idr-corollary page:197 type:corollary
\par\Needspace{8\baselineskip}
\begin{Corollary}[\BookCite{48}]\label{cor:shifted-idr-common-subspaces}
{Let $\bm r_0^{(i)}\ (i=1,\ldots,m)$ be collinear to each other, $\mathcal G_0^{(i)}=\mathcal K_n(A+\sigma_i I,\bm r_0^{(i)})$, $\mathcal S$ be a subspace of $\mathbb R^n$, and define sequences of subspaces $\mathcal G_j^{(i)}$ as}
\[\mathcal G_j^{(i)}=[I-\omega_j^{(i)}(A+\sigma_i I)](\mathcal G_{j-1}^{(i)}\cap\mathcal S),\quad\omega_j^{(i)}\ne0,\quad j=1,2,\ldots,\]
{then it holds that $\mathcal G_j^{(1)}=\mathcal G_j^{(2)}=\cdots=\mathcal G_j^{(m)}$.}
\end{Corollary}
% block: shifted-idr-collinear-approach page:197 type:ordinary-paragraph
{Corollary\ \ref{cor:shifted-idr-common-subspaces} implies that if we consider solving the following linear system (seed system) and shifted linear system:}
\[A\bm x=\bm b,\quad(A+\sigma I)\bm x^{\sigma}=\bm b,\]
{then the collinear approach, as described in the shifted BiCGSTAB method, $\bm r_i=\pi_i^{\sigma}\bm r_i^{\sigma}$ is promising, since $\bm r_i=\bm0$ and $\pi_i^{\sigma}\ne0$ lead to $\bm r_i^{\sigma}=\bm0$. Here $\bm r_i$ is the $i$th residual vector of the original linear system (seed system) $A\bm x=\bm b$ and $\bm r_i^{\sigma}$ is the $i$th residual vector of the shifted linear system $(A+\sigma I)\bm x^{\sigma}=\bm b$. We now describe how to compute $\bm r_{k+1}^{\sigma}$ from the information $\bm r_{k+1}$ when $\bm r_i=\pi_i^{\sigma}\bm r_i^{\sigma}\ (\pi_i^{\sigma}\in\mathbb C,\ i=0,1,\ldots,k)$.}
% block: shifted-idr-initial-residuals page:197 type:ordinary-paragraph
{From the initial step of the IDR($s$) method in Algorithm\ \ref{alg:idr}, the initial residual vectors of the seed system are given as}
\begin{equation}\label{eq:shifted-idr-seed-initial-residual}\bm r_{k+1}=\bm r_k-\omega_kA\bm r_k,\quad(k=0,\ldots,s-1).\end{equation}
{Similar to \eqref{eq:shifted-idr-seed-initial-residual}, the initial residual vectors of the shifted system are as follows:}
\begin{equation}\label{eq:shifted-idr-shifted-initial-residual}\bm r_{k+1}^{\sigma}=\bm r_k^{\sigma}-\omega_k^{\sigma}(A+\sigma I)\bm r_k^{\sigma}=(1-\sigma\omega_k^{\sigma})\bm r_k^{\sigma}-\omega_k^{\sigma}A\bm r_k^{\sigma},\end{equation}
{where the parameters $\omega_k^{\sigma}\in\mathbb C$ are unknown. Substituting the relations $\bm r_i=\pi_i^{\sigma}\bm r_i^{\sigma}\ (i=0,\ldots,k)$ into \eqref{eq:shifted-idr-seed-initial-residual} yields}
\begin{equation}\label{eq:shifted-idr-scaled-initial-residual}\bm r_{k+1}^{\sigma}=\frac{\pi_k^{\sigma}}{\pi_{k+1}^{\sigma}}\bm r_k^{\sigma}-\omega_k\frac{\pi_k^{\sigma}}{\pi_{k+1}^{\sigma}}A\bm r_k^{\sigma}.\end{equation}
% block: shifted-idr-initial-parameters page:198 type:ordinary-paragraph
{From the relations $\bm r_i=\pi_i^{\sigma}\bm r_i^{\sigma}$, \eqref{eq:shifted-idr-seed-initial-residual} and residual polynomial $R_i(0)=1$, $\bm r_i=R_i(A)\bm r_0$, it follows that}
\[\begin{cases}\pi_k^{\sigma}=R_k(-\sigma),\\\pi_{k+1}^{\sigma}=\pi_k^{\sigma}+\omega_k\sigma\pi_k^{\sigma}.\end{cases}\]
{Thus, parameter $\pi_{k+1}^{\sigma}$ in \eqref{eq:shifted-idr-scaled-initial-residual} can be obtained when $\pi_k^{\sigma}$ is computed. Comparing the coefficients of \eqref{eq:shifted-idr-shifted-initial-residual} and \eqref{eq:shifted-idr-scaled-initial-residual}, we obtain $\omega_k^{\sigma}=\omega_k\pi_k^{\sigma}/\pi_{k+1}^{\sigma}$. The corresponding approximate solution of the shifted system is written as}
\begin{equation}\label{eq:shifted-idr-initial-solution}\bm x_{k+1}^{\sigma}=\bm x_k^{\sigma}+\omega_k^{\sigma}\bm r_k^{\sigma}.\end{equation}
% block: shifted-idr-main-seed page:198 type:ordinary-paragraph
{Next, we describe the derivation process of $\bm r_{k+1}^{\sigma}$ by using $\bm r_{k+1}$ in the main step of the IDR($s$) method in Algorithm\ \ref{alg:idr}. From \eqref{eq:idr-first-projected-vector} and \eqref{eq:idr-first-reduced-residual}, $\bm v_k$ and $\bm r_{k+1}$ of the IDR($s$) method are written as}
\begin{equation}\label{eq:shifted-idr-seed-projected-vector}\bm v_k=\bm r_k-\sum_{l=1}^{s}\gamma_l\Delta\bm r_{k-l}=(1-\gamma_1)\bm r_k+\sum_{l=1}^{s-1}(\gamma_l-\gamma_{l+1})v\bm r_{k-l}+\gamma_s\bm r_{k-s},\end{equation}
\begin{equation}\label{eq:shifted-idr-seed-main-residual}\bm r_{k+1}=(I-\omega_jA)\bm v_k.\end{equation}
{It follows from \eqref{eq:shifted-idr-seed-projected-vector} and \eqref{eq:shifted-idr-seed-main-residual} that we obtain}
\begin{equation}\label{eq:shifted-idr-seed-expanded-residual}\bm r_{k+1}=(I-\omega_jA)\left((1-\gamma_1)\bm r_k+\sum_{l=1}^{s-1}(\gamma_l-\gamma_{l+1})\bm r_{k-l}+\gamma_s\bm r_{k-s}\right).\end{equation}
% block: shifted-idr-main-shifted-residual page:198 type:ordinary-paragraph
{Residual vectors of the shifted system in the same form of \eqref{eq:shifted-idr-seed-expanded-residual} are defined as follows:}
\begin{equation}\label{eq:shifted-idr-shifted-expanded-residual}\begin{aligned}\bm r_{k+1}^{\sigma}&=(I-\omega_j^{\sigma}(A+\sigma I))\left((1-\gamma_1^{\sigma})\bm r_k^{\sigma}+\sum_{l=1}^{s-1}(\gamma_l^{\sigma}-\gamma_{l+1}^{\sigma})\bm r_{k-l}^{\sigma}+\gamma_s^{\sigma}\bm r_{k-s}^{\sigma}\right)\\
&=\left(I-\frac{\omega_j^{\sigma}}{1-\sigma\omega_j^{\sigma}}A\right)(1-\sigma\omega_j^{\sigma})\left((1-\gamma_1^{\sigma})\bm r_k^{\sigma}+\sum_{l=1}^{s-1}(\gamma_l^{\sigma}-\gamma_{l+1}^{\sigma})\bm r_{k-l}^{\sigma}+\gamma_s^{\sigma}\bm r_{k-s}^{\sigma}\right),\end{aligned}\end{equation}
{where parameters $\omega_j^{\sigma},\gamma_1^{\sigma},\ldots,\gamma_s^{\sigma}\in\mathbb C$ are unknown. Substituting $\bm r_i=\pi_i^{\sigma}\bm r_i^{\sigma}\ (i=0,1,\ldots,k+1)$ into \eqref{eq:shifted-idr-seed-expanded-residual} yields}
\begin{equation}\label{eq:shifted-idr-scaled-main-residual}\bm r_{k+1}^{\sigma}=(I-\omega_jA)\left((1-\gamma_1)\frac{\pi_k^{\sigma}}{\pi_{k+1}^{\sigma}}\bm r_k^{\sigma}+\sum_{l=1}^{s-1}(\gamma_l-\gamma_{l+1})\frac{\pi_{k-l}^{\sigma}}{\pi_{k+1}^{\sigma}}\bm r_{k-l}^{\sigma}+\gamma_s\frac{\pi_{k-s}^{\sigma}}{\pi_{k+1}^{\sigma}}\bm r_{k-s}^{\sigma}\right).\end{equation}
% block: shifted-idr-main-parameters page:199 type:ordinary-paragraph
{It follows from \eqref{eq:shifted-idr-seed-expanded-residual}, $\bm r_i=\pi_i^{\sigma}\bm r_i^{\sigma}\ (i=0,1,\ldots,k+1)$, and $R_i(0)=1$ that we obtain $\pi_{k+1}^{\sigma}$ as follows:}
\begin{equation}\label{eq:shifted-idr-main-pi}\pi_{k+1}^{\sigma}=(1+\omega_j\sigma)\left((1-\gamma_1)\pi_k^{\sigma}+\sum_{l=1}^{s-1}(\gamma_l-\gamma_{l+1})\pi_{k-l}^{\sigma}+\gamma_s\pi_{k-s}^{\sigma}\right).\end{equation}
{Comparing the corresponding coefficients of \eqref{eq:shifted-idr-shifted-expanded-residual} and \eqref{eq:shifted-idr-scaled-main-residual} gives}
\begin{equation}\label{eq:shifted-idr-coefficient-matching}\begin{cases}\dfrac{\omega_j^{\sigma}}{1-\sigma\omega_j^{\sigma}}=\omega_j,\\
1-\gamma_1^{\sigma}=\dfrac{1-\gamma_1}{1-\sigma\omega_j^{\sigma}}\times\dfrac{\pi_k^{\sigma}}{\pi_{k+1}^{\sigma}},\\
\gamma_l^{\sigma}-\gamma_{l+1}^{\sigma}=\dfrac{\gamma_l-\gamma_{l+1}}{1-\sigma\omega_j^{\sigma}}\times\dfrac{\pi_{k-l}^{\sigma}}{\pi_{k+1}^{\sigma}},\\
\gamma_s^{\sigma}=\dfrac{\gamma_s}{1-\sigma\omega_j^{\sigma}}\times\dfrac{\pi_{k-s}^{\sigma}}{\pi_{k+1}^{\sigma}}\end{cases}\end{equation}
{for $l=1,2,\ldots,s-1$. From \eqref{eq:shifted-idr-coefficient-matching}, parameters $\omega_k^{\sigma},\gamma_1^{\sigma},\ldots,\gamma_s^{\sigma}$ are determined as}
\begin{equation}\label{eq:shifted-idr-parameters}\begin{cases}\omega_j^{\sigma}=\dfrac{\omega_j}{1+\sigma\omega_j},\\
\gamma_1^{\sigma}=1-(1-\gamma_1)(1+\sigma\omega_j)\dfrac{\pi_k^{\sigma}}{\pi_{k+1}^{\sigma}},\\
\gamma_{l+1}^{\sigma}=\gamma_l^{\sigma}-(\gamma_l-\gamma_{l+1})(1+\sigma\omega_j)\dfrac{\pi_{k-l}^{\sigma}}{\pi_{k+1}^{\sigma}}\end{cases}\end{equation}
{for $l=1,2,\ldots,s-1$.}
% block: shifted-idr-main-solution page:199 type:ordinary-paragraph
{Using the relation $\bm r_{k+1}^{\sigma}=\bm b-(A+\sigma I)\bm x_{k+1}^{\sigma}$, the approximate solutions $\bm x_{k+1}^{\sigma}$ can be derived from \eqref{eq:shifted-idr-shifted-expanded-residual} as}
\begin{equation}\label{eq:shifted-idr-main-solution}\bm x_{k+1}^{\sigma}=\bm x_k^{\sigma}+\omega_j^{\sigma}\bm v_k^{\sigma}-\sum_{l=1}^{s}\gamma_l^{\sigma}\Delta\bm x_{k-l}^{\sigma},\end{equation}
{where $\bm v_k^{\sigma}=\bm r_k^{\sigma}-\sum_{l=1}^{s}\gamma_l^{\sigma}\Delta\bm r_{k-l}^{\sigma}$. For practical computation, one can express $\bm v_k^{\sigma}$ by the residual vectors $\bm r_{k-s},\ldots,\bm r_k$ of the seed system. All the steps above give the shifted IDR($s$) method as described in Algorithm\ \ref{alg:shifted-idr}.}
% block: shifted-idr-algorithm page:200 type:ordinary-paragraph
\begin{algorithm}[H]\caption{{The shifted IDR($s$) method (Seed system: $(A+\sigma_f I)\bm x=\bm b$)}}\label{alg:shifted-idr}\begin{algorithmic}[1]
\Statex\textbf{Input:} $\bm x_0^{(i)}=\bm0,\ \bm r_0=\bm b,\ P\in\mathbb C^{n\times s},\ \pi_0^{(f,i)}=1,\ \sigma_i$ {for} $i=1,\ldots,m$
\Statex\textbf{Output:} $\bm x^{(i)}$ {for} $i=1,\ldots,m$
\For{$k=0,1,\ldots,s-1$}
\State $\bm v=(A+\sigma_f I)\bm r_k,\ \omega=\bm v^{\mathrm H}\bm r_k/\bm v^{\mathrm H}\bm v$
\State $\Delta X(:,k+1)=\omega\bm r_k,\ \Delta R(:,k+1)=-\omega\bm v$
\State $\bm x_{k+1}^{(f)}=\bm x_k^{(f)}+\Delta X(:,k+1),\ \bm r_{k+1}=\bm r_k+\Delta R(:,k+1)$
\State ({Iteration for shifted system})
\For{$i(\ne f)=1,2,\ldots,m$}
\If{$\|\bm r_k^{(i)}\|>\varepsilon\|\bm b\|$}
\State $\pi_{k+1}^{(f,i)}=\pi_k^{(f,i)}+\omega(\sigma_i-\sigma_f)\pi_k^{(f,i)},\quad\omega^{(f,i)}=\dfrac{\omega}{1+\omega(\sigma_i-\sigma_f)}$
\State $\bm x_{k+1}^{(i)}=\bm x_k^{(i)}+\dfrac{\omega^{(f,i)}}{\pi_k^{(f,i)}}\bm r_k$
\EndIf\EndFor\EndFor
\State $j=1,\ k=s,\ M=P^{\mathrm H}\Delta R,\ \bm h=P^{\mathrm H}\bm r_k$
\While{{stopping criterion is not satisfied}}
\For{$l=0,1,\ldots,s$}
\State {Solve $\bm c$ from $M\bm c=\bm h$}
\State $\bm q=-\Delta R\bm c,\ \bm v=\bm r_k+\bm q$
\If{$l=0$}
\State $\bm t=(A+\sigma_f I)\bm v,\ \omega=\bm t^{\mathrm H}\bm v/\bm t^{\mathrm H}\bm t$
\State $\Delta R(:,j)=\bm q-\omega\bm t,\ \Delta X(:,j)=-\Delta X\bm c+\omega\bm v$
\Else
\State $\Delta X(:,j)=-\Delta X\bm c+\omega\bm v,\ \Delta R(:,j)=-(A+\sigma_f I)\Delta X(:,j)$
\EndIf
\State $\bm r_{k+1}=\bm r_k+\Delta R(:,j),\ \bm x_{k+1}^{(f)}=\bm x_k^{(f)}+\Delta X(:,j)$
\State $\delta\bm m=P^{\mathrm H}\Delta R(:,j),\ M(:,j)=\delta\bm m,\ \bm h=\bm h+\delta\bm m$
\algstore{shiftedidrcontinued}
\end{algorithmic}\end{algorithm}
\begin{algorithm}[H]\ContinuedFloat\caption[]{{The shifted IDR($s$) method (continued)}}\begin{algorithmic}[1]
\algrestore{shiftedidrcontinued}
\State ({Iteration for shifted system})
\State $\gamma_1=c_{j-1},\ \gamma_2=c_{j-2},\ldots,\gamma_{j-1}=c_1$
\State $\gamma_j=c_s,\ \gamma_{j+2}=c_{s-1},\ldots,\gamma_s=c_j$
\For{$i(\ne f)=1,2,\ldots,m$}
\If{$\|\bm r_k^{(i)}\|>\varepsilon\|\bm b\|$}
\State $\alpha_i=1+\omega(\sigma_i-\sigma_f)$
\State $\pi_{k+1}^{(f,i)}=\alpha_i\left((1-\gamma_1)\pi_k^{(f,i)}+\sum_{g=1}^{s-1}(\gamma_g-\gamma_{g+1})\pi_{k-g}^{(f,i)}+\gamma_s\pi_{k-s}^{(f,i)}\right)$
\State $\gamma_1^{(f,i)}=1-\alpha_i(1-\gamma_1)\pi_k^{(f,i)}/\pi_{k+1}^{(f,i)}$
\State $\gamma_{g+1}^{(f,i)}=\gamma_g^{(f,i)}-(\gamma_g-\gamma_{g+1})\alpha_i\pi_{k-g}^{(f,i)}/\pi_{k+1}^{(f,i)}\quad(g=1,\ldots,s-1)$
\State $\bm x_{k+1}^{(i)}=\bm x_k^{(i)}+\omega\bm v/\pi_{k+1}^{(f,i)}-\sum_{g=1}^{s}\gamma_g^{(f,i)}(\bm x_{k+1-g}^{(i)}-\bm x_{k-g}^{(i)})$
\EndIf\EndFor
\State $k=k+1,\ j=j+1$
\State $j=(j-1)\%s+1$ {($\%$: modulo operation, i.e. $a\%n=r$, where $a=mn+r$.)}
\EndFor\EndWhile\end{algorithmic}\end{algorithm}
